\section{Introduction}\label{sec:intro}
All logarithms are natural. Write $\sfq$ for the positive squarefree integers,
$\mu$ for the M\"obius function, and $P^+(d)$ for the largest prime factor of
$d$, with $P^+(1)=1$. Put
\[
 c_0=\frac{6}{\pi^2},\qquad
 \nu_d=\frac1d\prod_{p\le P^+(d)}\frac{p}{p+1},\qquad
 \delta_d=c_0\nu_d\quad(d\in\sfq).
\]
Thus $\nu_1=1$. The primes in each product are unrestricted by divisibility.

The squarefree integers form a rooted tree: the root is $1$, and the parent
of $d>1$ is $d/P^+(d)$. Write $j\preceq k$ when $j$ is an ancestor of $k$,
including equality. Equivalently, $j\mid k$ and every prime divisor of $k/j$
exceeds $P^+(j)$. Listing the prime factors of $k$ in increasing order, the
ancestors of $k$ are exactly the products of its initial segments.

Define a real symmetric matrix on $\N\times\N$ by
\begin{equation}\label{eq:K}
 K_{jk}=
 \begin{cases}
 \mu(k/j)\delta_k,&j,k\in\sfq,\ j\preceq k,\\
 \mu(j/k)\delta_j,&j,k\in\sfq,\ k\preceq j,\\
 0,&\text{otherwise}.
 \end{cases}
\end{equation}
The two applicable cases agree on the diagonal. Nonsquarefree coordinates
are zero. We shall show that this matrix defines a positive compact operator
on $\ell^2(\N)$.

For a positive compact operator $A$, let
\[
 \ncount(A;t)=\#\{r:\lambda_r(A)>t\},\qquad t>0,
\]
where positive eigenvalues are ordered decreasingly, with multiplicity.
All counting functions in the operator arguments use this strict convention.
The subject of this paper is the rate at which the eigenvalues of $K$ tend
to zero.

\begin{theorem}\label{thm:main}
The operator $K$ in \eqref{eq:K} is positive and compact. As $s\downarrow0$,
\begin{equation}\label{eq:maincount}
 \ncount(K;s)\sim\frac{c_0}{s\log(1/s)}.
\end{equation}
Equivalently, as $r\to\infty$,
\begin{equation}\label{eq:maineigen}
 \lambda_r(K)\sim\frac{c_0}{r\log r}.
\end{equation}
Consequently,
\[
 \sum_{r\le R}\lambda_r(K)\sim c_0\log\log R.
\]
The same three asymptotics hold for the positive compact operator
\begin{equation}\label{eq:deflation}
 L=K-\frac{(Ke_1)(Ke_1)^*}{c_0}.
\end{equation}
\end{theorem}

Two immediate consequences are that $K$ and $L$ lie in the Schatten class
$\mathcal S_p$ exactly when $p>1$ (Corollary~\ref{cor:schatten}), and that
$K$ is not trace class even though $K$ minus its diagonal is
(Proposition~\ref{prop:traceclass}).

\subsection{Where the kernel comes from}
The kernel $K$ appeared in the author's earlier study~\cite{Kline} of a
sparse matrix, introduced in~\cite{Kline19}, whose determinant is the
Mertens function $M(n)=\sum_{m\le n}\mu(m)$. Let $A_n$ be
the $n\times n$ matrix with ones on the diagonal and an additional one in
position $(i,i/P^+(i))$ for each squarefree $i$ with $1<i\le n$. Its inverse
$C_n=A_n^{-1}$ has entries $\mu(i/j)$ when $j\preceq i$
(Section~\ref{sec:finite}). Replacing the first row of $A_n$ by ones gives
a matrix $B_n$ with $\det B_n=M(n)$. The following facts are proved
in~\cite{Kline}.
\begin{itemize}
\item \cite[Theorem~4.3]{Kline}: the matrices $C_n^*C_n/n$, extended by zero,
 converge to $K$ in Hilbert--Schmidt norm. Hence, for every fixed $r$,
 $s_r(C_n)/\sqrt n\to\sqrt{\lambda_r(K)}$.
\item \cite[Theorem~4.6]{Kline}: for every fixed $r\ge1$,
 $\sqrt n\,\sigma_{n-r}(B_n)\to\lambda_r(L)^{-1/2}$, with $L$ as in
 \eqref{eq:deflation}.
\end{itemize}
Both statements fix the index $r$ and let $n\to\infty$. They identify each
limit but do not describe how the limits behave as $r$ grows. The earlier
analysis also gives little quantitative information about the limiting
operator itself: Hilbert--Schmidt membership gives
$\sum_r\lambda_r(K)^2<\infty$, but not the order of $\lambda_r(K)$.
Theorem~\ref{thm:main} answers this question for the limiting operators.
It does not, by itself, give information about $s_{r}(C_n)$ when $r$ and
$n$ grow together; Section~\ref{sec:finite} records what an operator-norm
error estimate would add.

A first guess comes from the diagonal. The diagonal entries of $K$ are
$\delta_d$, and Mertens' product theorem gives
$\delta_d\approx e^{-\gamma}/(d\log P^+(d))$ when $P^+(d)$ is large. The
off-diagonal part is not small in the sense needed, however. It is trace
class (Proposition~\ref{prop:traceclass}), which controls partial sums of
eigenvalues up to a bounded error. The eigenvalues themselves, however,
are of order $1/(r\log r)$, and a trace-class operator may have singular
values that are not $o(1/(r\log r))$. Trace-class membership alone
therefore does not transfer an asymptotic for individual eigenvalues from
the diagonal to $K$. The proof uses the structure of $K$ instead.

\subsection{Outline of the proof}
The argument has four steps.

\emph{Removing the signs.} On comparable squarefree indices
$\mu(k/j)=\mu(j)\mu(k)$, so conjugation by the diagonal unitary with
entries $\mu(d)$ replaces $K$ by the unsigned kernel with entries
$\delta_{\max(j,k)}$ on comparable pairs. The spectrum is unchanged.

\emph{A cylinder representation.} Choose a random set of primes, including
each prime $p$ independently with probability $1/(p+1)$. For $d\in\sfq$ let
$A_d$ be the event that the selected primes up to $P^+(d)$ are exactly the
prime divisors of $d$. Then $A_d$ has probability $\nu_d$; the events are
nested along the tree and disjoint for incomparable indices. The unsigned
kernel is therefore $c_0$ times the Gram matrix of the indicators
$\one_{A_d}$ (Section~\ref{sec:kernel}).

\emph{From masses to eigenvalues.} Section~\ref{sec:tree} treats any tree
of nested sets whose children are disjoint and carry at most a fixed
fraction $\theta<1$ of their parent's mass. The operator $T$, the sum
over vertices $v$ of the rank-one maps
$f\mapsto\inner{f}{\one_{A_v}}\one_{A_v}$, is compact, and splitting each subtree into
its root and its children expresses $T$ as a rank-one operator plus an
orthogonal sum of rescaled copies of the subtree operators. Rank-one
interlacing, applied recursively, traps the eigenvalue count of $T$
between the number of leaves and an inflated number of vertices of the
finite tree $\{v:\nu(A_v)>t\}$. When the child-to-parent mass ratios tend
to zero outside finite sets and the mass count is regularly varying, both
bounds are asymptotic to the mass count (Theorem~\ref{thm:transfer}).

\emph{Counting the masses.} It remains to count $d$ with $\nu_d>t$
(Section~\ref{sec:arithmetic}). Grouping $d$ by its largest prime factor
reduces this to counts of squarefree smooth numbers, which we evaluate
with the Dickman function and Mertens' theorems. The answer is
$D(t)\sim1/(t\log(1/t))$. Its constant is the product of the squarefree
density $c_0$, the limit $\zeta(2)e^{-\gamma}$ of $m(p)\log p$, where
$m(p)=\prod_{q\le p}q/(q+1)$, and the classical value
$\int_0^\infty\rho(v)\,dv=e^\gamma$; since $c_0\zeta(2)=1$, the product
is one. The factor $c_0$ in Theorem~\ref{thm:main} comes only from the
relation $\delta_d=c_0\nu_d$.

\subsection{Scope}
The result is unconditional. Its ingredients are classical: rank-one
interlacing, the min--max principle, smooth-number estimates
from~\cite{HT}, and Mertens' theorems. What this paper contributes is the
cylinder representation of $K$, the reduction of its eigenvalue count to a
count of cylinder masses, and the evaluation of that count, which together
give the exact asymptotic in Theorem~\ref{thm:main}. Section~\ref{sec:related}
discusses related results on arithmetic matrices, operators on trees, and
self-similar spectral problems, and lists comparisons that remain open.

We do not obtain a rate of convergence or a second-order term in
\eqref{eq:maincount}. The theorem concerns the limiting operators; it
does not give a growing-index statement for the finite matrices without
an additional estimate (Corollary~\ref{cor:growing}). The kernel is
positive after a diagonal change of signs, and the proof uses no
cancellation in sums of $\mu$; in particular, no new estimate for $M(n)$
follows.

Appendix~\ref{sec:profiles} gives a second sufficient condition for the
transfer from masses to eigenvalues, based on an exact rank-one shift
identity, and Appendix~\ref{sec:stationary} applies it to stationary
trees and computes the spectrum of homogeneous trees. Neither appendix is
needed for Theorem~\ref{thm:main}. Whether the transfer requires an
assumption beyond regular variation of the mass count is left open
(Question~\ref{q:drop}).
